\begin{sol}[Péndulo suspendido de un carro móvil]

\subsubsection*{Pregunta 1: Grados de libertad}

El sistema tiene $2$ grados de libertad: la posición del carro a lo largo del riel y la orientación del péndulo. Las coordenadas generalizadas naturales son $q_1 = x$ (posición horizontal del carro) y $q_2 = \theta$ (ángulo de la varilla respecto a la vertical).

\subsubsection*{Pregunta 2: Lagrangiano}

Posición de la masa pendular respecto a un origen fijo, con $O$ (el pivote en el carro) a altura de referencia $y=0$:
$$X = x + \ell\sin\theta, \qquad Y = -\ell\cos\theta$$

Velocidades:
$$\dot{X} = \dot{x} + \ell\dot\theta\cos\theta, \qquad \dot{Y} = \ell\dot\theta\sin\theta$$

Energía cinética (carro + partícula):
$$T = \frac{1}{2}M\dot{x}^2 + \frac{1}{2}m\left(\dot{X}^2+\dot{Y}^2\right)$$

Expandiendo $\dot{X}^2+\dot{Y}^2 = \dot{x}^2 + 2\ell\dot{x}\dot\theta\cos\theta + \ell^2\dot\theta^2$:

$$T = \frac{1}{2}(M+m)\dot{x}^2 + m\ell\dot{x}\dot\theta\cos\theta + \frac{1}{2}m\ell^2\dot\theta^2$$

Energía potencial (el carro no cambia de altura):
$$V = mgY = -mg\ell\cos\theta$$

$$\boxed{L = \frac{1}{2}(M+m)\dot{x}^2 + m\ell\dot{x}\dot\theta\cos\theta + \frac{1}{2}m\ell^2\dot\theta^2 + mg\ell\cos\theta}$$

\subsubsection*{Pregunta 3: Ecuaciones de movimiento}

\textbf{Respecto a $x$:}
$$\frac{\partial L}{\partial\dot{x}} = (M+m)\dot{x} + m\ell\dot\theta\cos\theta, \qquad \frac{\partial L}{\partial x} = 0$$

$$\boxed{(M+m)\ddot{x} + m\ell\left(\ddot\theta\cos\theta - \dot\theta^2\sin\theta\right) = 0}$$

\textbf{Respecto a $\theta$:}
$$\frac{\partial L}{\partial\dot\theta} = m\ell\dot{x}\cos\theta + m\ell^2\dot\theta, \qquad \frac{\partial L}{\partial\theta} = -m\ell\dot{x}\dot\theta\sin\theta - mg\ell\sin\theta$$

Al aplicar Euler-Lagrange, los términos $\mp m\ell\dot{x}\dot\theta\sin\theta$ se cancelan, quedando:

$$\boxed{\ell\ddot\theta + \ddot{x}\cos\theta + g\sin\theta = 0}$$

Estas dos ecuaciones están acopladas: la aceleración del carro depende de $\theta$ y $\dot\theta$, y la aceleración angular depende de $\ddot{x}$.

\subsubsection*{Pregunta 4: Coordenada cíclica y cantidad conservada}

El Lagrangiano no depende explícitamente de $x$ (solo aparece $\dot{x}$), por lo que $x$ es una \textbf{coordenada cíclica}. Su momento canónico conjugado se conserva:

$$p_x = \frac{\partial L}{\partial\dot{x}} = (M+m)\dot{x} + m\ell\dot\theta\cos\theta = \text{constante}$$

\paragraph{Interpretación física} $p_x$ es exactamente el momento lineal horizontal total del sistema (carro más partícula): $p_x = M\dot x + m\dot X$. Su conservación refleja que no existe ninguna fuerza externa horizontal actuando sobre el sistema completo (el riel es sin fricción y la gravedad es vertical), consistente con el principio de conservación de momento lineal en ausencia de fuerzas externas netas en esa dirección.

\end{sol}
